% Chapter 5: Meter Dynamics with Non-Commuting Hamiltonians
\chapter{Quantifying Non-Commuting Meter Dynamics}\label{Ch4:NonCommutingMeter}

\qquad Chapter~\ref{Ch3:SpeedingUpQuantumAnnealing} established that the QND-like meter protocol achieves a linear speedup factor of $(1+m_{\max})$ when two conditions hold: the meter observable $X_M$ commutes with the meter Hamiltonian $H_M$, and the meter is initialized in an eigenstate of $X_M$. Under these conditions the meter remains frozen in its initial state, and the system evolves under a single, time-independent rescaling of its Hamiltonian. A natural question is what happens when this commutation condition is relaxed.

\qquad Any physical auxiliary system will possess internal dynamics that do not commute perfectly with the coupling observable. When $[H_M, X_M] \neq 0$, three effects arise: (i) the meter undergoes coherent precession away from its initial state, (ii) the system becomes entangled with the meter, and (iii) the effective energy rescaling becomes time-dependent rather than constant. We derive an analytical from for these effects for a qubit meter, and validate it its predictions against exact numerical simulations.

\begin{figure}[htbp]
    \centering
    \begin{subfigure}[t]{0.45\textwidth}
        \centering
        \refstepcounter{subfigure}\makebox[\textwidth][l]{\hspace{-2mm}(\thesubfigure)}\label{fig:IonQDetails_a}
        \vspace{-1.2ex}
            \includegraphics[width=0.6\textwidth]{connectivity.png}
    \end{subfigure}%
    \hfill
    \begin{subfigure}[t]{0.55\textwidth}
        \centering
        \refstepcounter{subfigure}\makebox[\textwidth][l]{\hspace{-2mm}(\thesubfigure)}\label{fig:IonQDetails_b}
        \vspace{-1.2ex}
        \includegraphics[width=0.95\textwidth]{Oscillator.png}
    \end{subfigure}
    \caption[IOQ Architecture]{(a) All-to-all connectivity of IonQ quantum computer cartoon, makes it easy to create $X_M \otimes H_S(t)$. (b) Ion trapped in external field. \textit{Source: IonQ Media Resources (\url{https://ionq.com})}}
    \label{fig:IonQDetails}
\end{figure}

\qquad This is physically motivated by annealing implementations on ion-trap devices where we can achieve a global coupling of the meter to the system Hamiltonian but more likely at the cost of the meter having its own intrinsic dynamics. Ion trap architecture such as those offered by IonQ feature this high qubit connectivity allowing for a qubit to be used as a meter with some global coupling via an electromagnetic wave\cite{IONQArichitecture2019}, at the cost of the qubit having its own intrinsic dynamics, as it is subject to external fields (e.g., laser or microwave fields) that drive its internal transitions.

%----------------------------------------------------------------------------------
\section{Joint Hamiltonian and the Role of Non-Commutation}\label{sec:joint_hamiltonian}
%----------------------------------------------------------------------------------

\qquad We try to understand this behaviour by writing the full system--meter Hamiltonian from Chapter~\ref{Ch3:SpeedingUpQuantumAnnealing} in its most general form,
%%
\begin{equation}\label{eq:total_hamiltonian_nc}
    H(t) = H_S(t) \otimes \mathds{1}_M + x_0\, H_S(t) \otimes X_M + \mathds{1}_S \otimes H_M,
\end{equation}
%%
where $H_S(t)$ is the time-dependent system Hamiltonian, $X_M$ is the meter observable that couples to the system, $x_0$ is the coupling strength, and $H_M$ governs the free meter dynamics as before. Throughout this chapter we specialise to a qubit meter with $X_M = \sigma_z, H_M = \omega\,\sigma_x$ dynamics so that $[H_M, X_M] = [\omega\sigma_x, \sigma_z] = -2i\omega\sigma_y \neq 0$ for any $\omega \neq 0$. $\omega$ allows us to quantify the difference from an ideal commuting regime. The total Hamiltonian becomes \color{red} Shorten this \color{black}
%%
\begin{equation}\label{eq:total_hamiltonian_explicit}
    H_{\text{tot}}(t) = H_S(t) \otimes \mathds{1}_M + x_0\, H_S(t) \otimes \sigma_z + \mathds{1}_S \otimes \omega\,\sigma_x.
\end{equation}
%%
In the commuting case ($\omega = 0$), the Hamiltonian is diagonal in the $\sigma_z$ eigenbasis of the meter, From Ch. \ref{Ch3:SpeedingUpQuantumAnnealing} we know the evolution splits into two non-interacting branches, each with a constant rescaling factor $(1 \pm x_0)$. The meter, once initialized in $|0\rangle$, remains there for all time. In contrast, when $\omega \neq 0$, the off-diagonal $\omega\sigma_x$ term mixes $|0\rangle$ and $|1\rangle$ at every instant, continuously coupling the two branches. The Hamiltonian is no longer block-diagonal in the meter basis, and no time-independent basis disentangles the system from the meter.

%----------------------------------------------------------------------------------
\section{Meter Dynamics in Each Energy Sector}\label{sec:meter_dynamics}
%----------------------------------------------------------------------------------

\qquad To develop an analytical picture, we work in the instantaneous eigenbasis of the system Hamiltonian, $H_S(t) = \sum_n E_n(t)\,|n(t)\rangle\langle n(t)|$. Within each energy sector labelled by $n$, the system eigenvalue $E_n(t)$ acts as a parameter, and the meter sees an effective $2 \times 2$ Hamiltonian. We derive this sector Hamiltonian in three steps. We act with the projector $|n(t)\rangle\langle n(t)| \otimes \mathds{1}_M$ on Eq.~\eqref{eq:total_hamiltonian_explicit}. Using $H_S(t)|n(t)\rangle = E_n(t)|n(t)\rangle$. This yields the effective meter Hamiltonian in sector $n$
%%
\begin{equation}\label{eq:meter_hamiltonian_sector}
    H_{\text{meter}}^{(n)}(t) = E_n(t)\,\mathds{1}_M + x_0\,E_n(t)\,\sigma_z + \omega\,\sigma_x.
\end{equation}
%%
\qquad The identity term $E_n(t)\,\mathds{1}_M$ produces only an overall phase and does not affect the meter's internal dynamics. The non-trivial part is $\tilde{H}_{\text{meter}}^{(n)}(t) = x_0\,E_n(t)\,\sigma_z + \omega\,\sigma_x$. In the basis $\{|0\rangle, |1\rangle\}$ (eigenstates of $\sigma_z$ with eigenvalues $\pm 1$), this matrix reads
%%
\begin{equation}\label{eq:meter_matrix}
    \tilde{H}_{\text{meter}}^{(n)}(t) = \begin{pmatrix} +x_0\,E_n(t) & \omega \\ \omega & -x_0\,E_n(t) \end{pmatrix}.
\end{equation}
%%
\qquad The diagonal entries $\pm x_0 E_n(t)$ are the energy splittings produced by the system meter coupling, and the off-diagonal entry $\omega$ is the transverse field that drives transitions between the meter states. This is precisely the Hamiltonian of a spin $\tfrac{1}{2}$ particle in a time-dependent longitudinal field $x_0 E_n(t)$ and a constant transverse field $\omega$. The eigenvalues and normalized eigenstates given as
%%
\begin{align}\label{eq:meter_eigenvalues}
    \lambda_{n,\pm}(t) = \pm\,\Omega_n(t) &= \pm \sqrt{\big(x_0\,E_n(t)\big)^2 + \omega^2} \nonumber \\
    |+_n(t)\rangle = \cos\theta_n(t)\,|0\rangle + \sin\theta_n(t)\,|1\rangle, \qquad
    &|-_n(t)\rangle = -\sin\theta_n(t)\,|0\rangle + \cos\theta_n(t)\,|1\rangle,
\end{align}
%%
where the mixing angle $\theta_n(t)$ satisfies
\begin{equation}\label{eq:mixing_angle}
    \tan\bigl(2\theta_n(t)\bigr) = \frac{\omega}{x_0\,E_n(t)}.
\end{equation}
%%
\qquad As before we initialize the meter in $|0\rangle$, the $+1$ eigenstate of $\sigma_z$ for a positive energy rescaling. Within sector $n$, the meter evolves under $\tilde{H}_{\text{meter}}^{(n)}(t)$. For simplicity, we make an adiabatic approximation such that $E_n(t)$ where it varies slowly compared to the Rabi frequency $\Omega_n(t)$. This is justified when the system's annealing timescale $T$ is much longer than $1/\Omega_n$. The initial meter state $|0\rangle$ is decomposed in the instantaneous eigenbasis of $\tilde{H}_{\text{meter}}^{(n)}(t)$ at $t=0$:
%%
\begin{equation}\label{eq:meter_decomposition}
    |0\rangle = \cos\theta_n(0)\,|+_n(0)\rangle - \sin\theta_n(0)\,|-_n(0)\rangle.
\end{equation}
%%
This follows directly from inverting Eq.~\eqref{eq:meter_eigenvalues}. Under the adiabatic-meter approximation, each component tracks the instantaneous eigenstate and acquires a dynamical phase under the eigenvalues $\pm\Omega_n(t)$:
%%
\begin{equation}\label{eq:meter_time_evolution}
    |\chi_n(t)\rangle = \cos\theta_n(0)\,e^{-i\Phi_n^{+}(t)}\,|+_n(t)\rangle - \sin\theta_n(0)\,e^{-i\Phi_n^{-}(t)}\,|-_n(t)\rangle,
\end{equation}
%%
\qquad The quantity that controls the effective system rescaling is the meter polarisation $\langle\sigma_z\rangle_n(t) = \langle\chi_n(t)|\sigma_z|\chi_n(t)\rangle$. A crucial point is that the amplitudes in Eq.~\eqref{eq:meter_time_evolution} are set by the initial mixing angle $\theta_n(0)$, while the matrix elements of $\sigma_z$ in the instantaneous eigenbasis depend on $\theta_n(t)$.  yields the central result for the meter dynamics:
%%
\begin{equation}\label{eq:sigma_z_final}
    \langle\sigma_z\rangle_n(t) = \frac{x_0^2\,E_n(0)\,E_n(t)}{\Omega_n(0)\,\Omega_n(t)} + \frac{\omega^2}{\Omega_n(0)\,\Omega_n(t)}\,\cos\!\left(2\int_0^t \Omega_n(s)\,ds\right),
\end{equation}
%%
with $\Omega_n(t) = \sqrt{(x_0 E_n(t))^2 + \omega^2}$. The first term depends on both the initial energy $E_n(0)$ and the instantaneous energy $E_n(t)$. The second term is an oscillatory component at twice the Rabi frequency. Its amplitude $\omega^2/(\Omega_n(0)\Omega_n(t))$ is largest when $\omega$ dominates both the initial and instantaneous coupling energies, and vanishes in the commuting limit $\omega \to 0$.

\begin{figure}[htbp]
    \centering
    \begin{subfigure}[t]{0.48\textwidth}
        \centering
        \refstepcounter{subfigure}\makebox[\textwidth][l]{\hspace{-2mm}(\thesubfigure)}\label{fig:NonCommutationRescaling_a}
        \vspace{-1.2ex}
            \includegraphics[width=0.95\textwidth]{NonCommutation/SigmaZTimeTrace.png}
    \end{subfigure}%
    \hfill
    \begin{subfigure}[t]{0.48\textwidth}
        \centering
        \refstepcounter{subfigure}\makebox[\textwidth][l]{\hspace{-2mm}(\thesubfigure)}\label{fig:NonCommutationRescaling_b}
        \vspace{-1.2ex}
        \includegraphics[width=0.95\textwidth]{NonCommutation/QARescale.png}
    \end{subfigure}
    \caption[Meter expectation value]{$\langle\sigma_z\rangle(t)$ for the Landau-Zener model and Ising model}
    \label{fig:NonCommutationRescaling}
\end{figure}

\qquad The total effect of the non-commuting dynamics on the speedup are found by time-averaging the meter polarisation in Eq.~\eqref{eq:sigma_z_final}. Over timescales much longer than the Rabi period $2\pi/\Omega_n$, the oscillatory cosine term averages to zero, leaving only the static component
%%
\begin{equation}\label{eq:sigma_z_time_avg}
    \overline{\langle\sigma_z\rangle}_n = \frac{x_0^2\,E_n(0)\,E_n(t)}{\Omega_n(0)\,\Omega_n(t)} = \frac{x_0^2\,E_n(0)\,E_n(t)}{\sqrt{(x_0 E_n(0))^2 + \omega^2}\,\sqrt{(x_0 E_n(t))^2 + \omega^2}}.
\end{equation}
%%
For the ground state ($n=0$), the effective rescaling factor becomes $1 + x_0\,\overline{\langle\sigma_z\rangle}_0$ rather than the ideal $(1+x_0)$. Since $|\overline{\langle\sigma_z\rangle}_0| \leq 1$ for all $\omega > 0$, with equality only at $\omega = 0$, the time-averaged rescaling is strictly reduced by non-commuting meter dynamics. Agreeing with our numerical simulations. The bound $1 + m_{\max}$ derived in Chapter~\ref{Ch3:SpeedingUpQuantumAnnealing} for the commuting case is therefore the global maximum, and any non-commuting meter dynamics strictly degrades the speedup. The degree of degradation is controlled by the ratio $\omega/(x_0 E_n)$: when $\omega \ll x_0 E_n$, the mixing angle $\theta_n \approx 0$ and the polarisation remains close to one, when $\omega \gg x_0 E_n$, the polarisation falls toward zero and the meter coupling becomes ineffective.

\color{red} TTS of the 1+max speedup \color{black}

%%==========================================================================
\section{Oscillatory Ripples}\label{sec:interference}
%%==========================================================================

\qquad The time-averaged term captures the effect of non-commuting dynamics but neglects the oscillatory term in Eq.~\eqref{eq:sigma_z_final}. At finite annealing times, the Rabi oscillations are not fully averaged out, and their signature appears as ripple modulations in the fidelity landscape. This section derives an analytical expression for this interference effect, between two branches. When the meter has non-commuting dynamics $H_M = \omega\sigma_x$, the relevant Rabi frequency for sector~$n$ is $\Omega_n(t) = \sqrt{\bigl(x_0\,E_n(t)\bigr)^2 + \omega^2}$. From Eq.~\eqref{eq:sigma_z_final}, the oscillatory part of $\langle\sigma_z\rangle_n$ defines the interference factor
%%
\begin{equation}\label{eq:interference_factor_general}
    C(T, \omega) = \frac{\omega^2}{\Omega_n(0)\,\Omega_n(T)}\,\cos\!\left(2\int_0^T \Omega_n(s)\,ds\right),
\end{equation}
%%
where the factor of 2 is the standard Rabi oscillation phase $2\int\Omega\,ds$ from the rotating-frame evolution. The amplitude prefactor $\omega^2/(\Omega_n(0)\,\Omega_n(T))$ is largest when $\omega$ dominates the coupling energy $x_0 E_n$ and vanishes in the commuting limit $\omega \to 0$. The zeros of the cosine, occurring when $2\int_0^T \Omega_n(s)\,ds = (2k+1)\pi/2$ for integer~$k$, define contours in the $(T, \omega)$ plane along which the oscillatory contribution vanishes.

\qquad We now compare the analytical predictions derived above against exact numerical simulations of the full system meter Schr\"odinger equation. The simulations use two test systems from Chapter~\ref{Ch3:SpeedingUpQuantumAnnealing}: (i) the Landau-Zener qubit (Eq.~\ref{LandauZenerEquation}) and (ii) a 3-qubit random Ising all to all system (Eq.~\ref{IsingProblemEquation}). In both cases the meter is a qubit with $X_M = \sigma_z$ and $H_M = \omega\sigma_x$, initialised in $|0\rangle$, with coupling $x_0 = 1$. \color{red} Update figures \color{black}



\qquad Fig.~\ref{fig:NonCommutationRescaling} displays the deviation of the effective rescaling factor $1 + x_0\langle\sigma_z\rangle(t)$ from its ideal value $(1+x_0) = 2$ as a function of annealing time $T$ and meter frequency $\omega$. 

\begin{figure}[htbp]
    \centering
    \begin{subfigure}[t]{0.32\textwidth}
        \centering
        \refstepcounter{subfigure}\makebox[\textwidth][l]{\hspace{-2mm}(\thesubfigure)}\label{fig:interference_analysis_a}
        \vspace{-1.2ex}
        \includegraphics[width=\textwidth]{NonCommutation/WNonFidDiff.png}
    \end{subfigure}%
    \hfill
    \begin{subfigure}[t]{0.32\textwidth}
        \centering
        \refstepcounter{subfigure}\makebox[\textwidth][l]{\hspace{-2mm}(\thesubfigure)}\label{fig:interference_analysis_b}
        \vspace{-1.2ex}
        \includegraphics[width=\textwidth]{NonCommutation/InterferenceFactor.png}
    \end{subfigure}%
    \hfill
    \begin{subfigure}[t]{0.32\textwidth}
        \centering
        \refstepcounter{subfigure}\makebox[\textwidth][l]{\hspace{-2mm}(\thesubfigure)}\label{fig:interference_analysis_c}
        \vspace{-1.2ex}
        \includegraphics[width=\textwidth]{NonCommutation/WNonCommuting.png}
    \end{subfigure}
    \caption[Interference effects in non-commuting meter regime]{Interference effects in the non-commuting meter regime for the Landau-Zener model with $x_0 = 1$. (a) Numerical fidelity from full Schr\"odinger evolution. (b) Analytical fidelity prediction using the adiabatic $\langle\sigma_z\rangle$ formula (Eq.~\eqref{eq:sigma_z_final}) to compute the effective rescaling factor. (c) Residual $F_{\mathrm{num}} - F_{\mathrm{anal}}$, showing that deviations concentrate at large $\omega$ and $T$ where the meter adiabaticity condition breaks down.}
    \label{fig:interference_analysis}
\end{figure}

\qquad Fig.~\ref{fig:interference_analysis} provides a direct test of the interference factor derived in Section~\ref{sec:interference}. Three panels are presented: (a)~the numerical fidelity map, (b)~the analytical fidelity prediction from the adiabatic $\langle\sigma_z\rangle$ formula (Eq.~\eqref{eq:sigma_z_final}), and (c)~the residual $F_{\mathrm{num}} - F_{\mathrm{anal}}$. The analytical prediction captures the dominant structure---the monotonic fidelity decrease with $\omega$ and the oscillatory ripples at small $\omega$. The ripple contours, defined by $2\int_0^T \Omega_n(s)\,ds = (2k+1)\pi/2$, align with those observed numerically. The residual panel shows that the adiabatic approximation is accurate at small $\omega$ and moderate $T$, but systematic deviations appear at large $\omega$ and $T$ where the meter adiabaticity condition $|d\Omega_n/dt|/\Omega_n^2 \ll 1$ breaks down and system--meter entanglement effects beyond the sector decomposition become significant.

\qquad The physical interpretation is as follows. At small $\omega$, the Rabi oscillations of the meter are slow enough that they are not averaged out over the annealing time $T$. The cosine term in Eq.~\eqref{eq:sigma_z_final} alternates between $+1$ and $-1$, periodically enhancing and suppressing the effective dephasing. As $\omega$ increases, the oscillation frequency grows and the oscillatory contribution averages away, leaving only the monotonic suppression from the time-averaged term. The transition between these regimes is governed by the condition $\Omega_n T \sim 1$.